\section{Synchronization, queues, and lifetime}
\label{sec:sync}

\subsection{Queues in use}

The engine has a main graphics queue and, when available, a distinct
graphics-family queue for geometry/upload work.  The topology compute
pass is submitted on the graphics queue with a single-time command buffer
and a fence, because its descriptor set and staging buffers are serialized
by a global mutex.  The geometry pass writes the renderer's packed pools;
the cull, raster, shadow, and ray-tracing passes then read them on the same
queue or after the documented cross-queue wait.

\subsection{The barriers inside the topology submission}

The submission starts by clearing the device-local election buffer and then
barriers the fill against the compute stage:
\begin{lstlisting}[style=glsl,caption={Election-buffer clear barrier.},label={lst:fillbarrier}]
vkCmdFillBuffer(cmd, nodeWritten, 0, nodeCount * 4, 0);
b.srcStageMask  = TRANSFER;
b.srcAccessMask = TRANSFER_WRITE;
b.dstStageMask  = COMPUTE_SHADER;
b.dstAccessMask = SHADER_READ | SHADER_WRITE;
\end{lstlisting}
This is required because \code{atomicExchange(nodeWritten[node], 1)} elects
one vertex writer per node; without the barrier the dispatch could read a
stale flag from a previous chunk and skip a payload.

The submission ends with a compute$\rightarrow$host barrier:
\begin{lstlisting}[style=glsl,caption={Counter readback barrier.},label={lst:hostbarrier}]
VkMemoryBarrier2 b{};
b.srcStageMask  = VK_PIPELINE_STAGE_2_COMPUTE_SHADER_BIT;
b.srcAccessMask = VK_ACCESS_2_SHADER_WRITE_BIT;
b.dstStageMask  = VK_PIPELINE_STAGE_2_HOST_BIT;
b.dstAccessMask = VK_ACCESS_2_HOST_READ_BIT;
vkCmdPipelineBarrier2(cmd, &dep);
\end{lstlisting}
This makes the integer counters visible to the host-mapped memory.  The
vertex/index writes are not read by the host, so they do not need a host
barrier.  The geometry pass is followed by a pool barrier:
\begin{lstlisting}[style=glsl,caption={Pool visibility barrier.},label={lst:poolbarrier}]
b.srcStageMask  = COMPUTE_SHADER;
b.srcAccessMask = SHADER_WRITE;
b.dstStageMask  = VERTEX_INPUT | INDEX_INPUT | DRAW_INDIRECT |
                  COMPUTE_SHADER | ACCELERATION_STRUCTURE_BUILD;
b.dstAccessMask = VERTEX_ATTRIBUTE_READ | INDEX_READ |
                  SHADER_READ | ACCELERATION_STRUCTURE_READ;
\end{lstlisting}
This is deliberately broad: the same pool range feeds the raster draw, the
GPU cull, the shadow cascade, and the ray-tracing BLAS.  The barrier runs
once per chunk, not once per consumer.

\subsection{The fence and the deferred release}

The topology submission uses a fence.  The host waits on the fence before
reading the counters.  An overflow retry releases the unpublished
reservation immediately (no draw command can reference it) and clears the
chunk's level data before the provider reserves again.  A previously
published span is released through the staged free queue; the actual free
happens two cull frames later.  This is necessary because the previous
frame's cull or draw command buffers may still be in flight and still
reference that span.  The renderer does not wait for the whole device; it
tracks the cull frame slot.

\subsection{Consolidation}

The vegetation instance buffers and the SDF flame anchors follow a similar
pattern.  A per-chunk placement dispatch writes into a scratch ring; the
host reads the accepted instance count, allocates an exact persistent
buffer, and copies the scratch slice on the GPU.  The copy is submitted on
the main thread after the placement fence, so no queue submission happens
inside a deferred callback.  The consolidated instance buffer is
double-buffered; the active pair is read by the vegetation pass while the
inactive pair is filled by the copy.

\subsection{The queue-submission mutex}

All Vulkan submissions go through a single global mutex keyed by
\code{VkQueue}.  This is intentionally coarse: the engine exposes several
queue handles that can alias the same physical queue, and the microseconds
saved by per-queue locking are not worth a submission race.  The mutex
covers:
\begin{itemize}
\item frame submit and present;
\item \code{runSingleTimeCommands} and its async variant;
\item upload-manager jobs;
\item topology submissions;
\item the vegetation exact-copy submission.
\end{itemize}
\code{vkDeviceWaitIdle} also takes the mutex.  The earlier validation
failure ``\code{vkQueueSubmit2}: VkQueue is simultaneously used'' came from
calling \code{vkDeviceWaitIdle} without it while a worker was submitting.
The fix is in \code{VulkanApp::deviceWaitIdle}, and cleanup now stops the
CPU thread pools before the first idle wait.

\subsection{Deferred destruction order}

At shutdown the engine:
\begin{enumerate}
\item stops background threads (no new submissions);
\item waits for the device idle under the submission mutex;
\item drains pending command buffers;
\item runs deferred-destruction callbacks;
\item calls \code{clean()} on the app;
\item destroys tracked resources;
\item destroys the device.
\end{enumerate}
The order matters because deferred callbacks capture renderer state.  If
the renderer were destroyed first, a callback that accesses
\code{chunkBuffers} would read freed memory.  The cleanup comment in
\code{VulkanApp::cleanup} documents this exact hazard.

\subsection{Why no host-side mesh wait}

The host waits on the topology fence to read counters, not to read
vertices.  The wait is short (one dispatch), and it is per chunk on a
worker thread, not on the render thread.  The render thread only sees
already-published draw entries.  This is the difference between a
GPU-resident pipeline and a GPU-assisted CPU pipeline: the host still
synchronizes, but it synchronizes on integers, not on megabytes.
